%!TEX root = ../main.tex
\section{Analysis}
\label{sec:analysis}

Description quality and behavioral reliability, while positively correlated at the population level, diverge at precisely the points that matter for deployment decisions. GPT-5.2 ranks first on MQM (91.6) but fourth on active admittance (6\%), while Gemini ranks second on MQM (90.2) but first on admittance (90\%), resistance (0.91), and caption bias resistance (0.89). The Spearman correlations in Table~\ref{tab:cross-dim} are high overall ($\rho = 0.83$--$0.95$), yet they mask these consequential reversals because the bottom of the ranking is concordant across dimensions. Split-half reliability of $\rho = 0.979$ over 100 random splits (Table~\ref{tab:stability}) confirms that this divergence is stable, not an artifact of sample variance. The practical implication is clear: a benchmark that reports only MQM would rank GPT-5.2 and Gemini as near-equivalent, missing the fact that one fabricates answers for blurred elements 98\% of the time while the other admits uncertainty 90\% of the time.

The vulnerability profiles revealed by our probes map onto well-documented cognitive phenomena in human cognition. Presupposition embedding (inexist probes) is the most effective deception technique across all eight models. Llama~4, for example, resists explicit false values at 0.76 (contra) but drops to 0.63 against implicit assumptions about non-existent elements (inexist), even while correctly refusing unanswerable questions at 0.94. This asymmetry parallels findings from eyewitness testimony research, where definite articles (``the broken headlight'') reliably induce false memories of objects that were never present \citep{loftus1975leading}. Caption bias resistance reveals a separate and unexpected pattern. Phi-4 follows poisoned captions almost entirely ($R = 0.05$), yet Gemma resists at $R = 0.38$ despite scoring 7 MQM points lower (69.1 vs.\ 62.2). If caption dependency were simply a capability limitation, weaker models should be uniformly more susceptible. Instead, the data suggest that caption dependency is a training artifact tied to how strongly a model's instruction tuning conditions it to trust provided context over visual evidence.

The most striking behavioral asymmetry emerges in how models handle visual uncertainty across evaluation modes. In passive descriptions, GPT-5.2 acknowledges blurred elements 26\% of the time, but when directly questioned about those same elements, its admittance drops to 6\%. Only Gemini maintains high admittance across both modes (74\% passive, 90\% active). This pattern suggests that direct questions trigger a ``must answer'' bias, consistent with models trained under RLHF to maximize helpfulness \citep{sharma2023towards}. Open-ended descriptions provide more latitude for natural hedging (``the label appears obscured''), while targeted questions compress the response space toward definitive answers. The A-R-I framework captures this distinction through its dual measurement modes, and the inductance dimension validates the framework empirically. When blurred elements are inferable from context, models answer correctly 21--80\% of the time (GPT-5.2 at 74\%, Gemini at 80\%), but when elements are genuinely unrecoverable, correctness drops to 0--14\% ($n = 50$ per model). This gap confirms that the framework measures real reasoning rather than random fabrication.

These findings are robust to methodological choices. The probe designer ablation (Table~\ref{tab:ablation}) shows that switching from GPT-4o to Mistral Large~3 as the probe generator yields negligible differences in caption bias resistance ($\Delta = 0.002$, $p = 0.49$, Cliff's $\delta = -0.007$) and only a small effect on hallucination resistance ($\Delta = -0.057$, Cliff's $\delta = -0.16$), confirming that our results reflect genuine model behavior rather than idiosyncrasies of the probe generator. Scale validation further supports stability: resistance scores computed on 100 figures closely match those from the full 250-figure set, with maximum model-level deviation of 0.02 points (Table~\ref{tab:stability}).
